\section{Introduction}
\label{sec:intro}

% Recommended structure (1.5--2 columns of IEEEtran):
%   Paragraph 1: Problem context (industry / scaling / technology trend)
%   Paragraph 2: Why prior approaches fall short
%   Paragraph 3: This work in one paragraph, leading into contribution bullets
%   Bulleted contribution list (3--5 bullets)
%   Paragraph 4: Paper roadmap

\IEEEPARstart{T}{he} continuing scaling of integrated circuits and the
increasing complexity of modern System-on-Chip (SoC) designs make
[the specific EDA problem you address] a central challenge for
chip-design productivity and quality of results (\QoR). \textbf{[REPLACE
with one to two sentences of industry / EDA flow context, with a concrete
example.]}

Prior work on this problem can be broadly grouped into [N] families.
\rev{One family of methods \cite{PLACEHOLDER_family1_a, PLACEHOLDER_family1_b}
relies on [approach], which achieves [property] but suffers from
[limitation].} A second family \cite{PLACEHOLDER_family2_a, PLACEHOLDER_family2_b}
adopts [different approach], achieving [property] but at the cost of
[other limitation]. To our knowledge, no existing method achieves all
three of \textit{[property A]}, \textit{[property B]}, and \textit{[property C]}
simultaneously.

In this paper, we propose \method{}, a [class of artifact, \eg{} algorithm,
framework, methodology] for [specific EDA problem]. \method{} combines
[ingredient 1] with [ingredient 2] in a [novel way / formulation], and is
designed to achieve [property A], [property B], and [property C]
simultaneously. We evaluate \method{} on [benchmark suites] using the
[PDK, tool flow]. Experimental results show that \method{} achieves
[headline result 1] and [headline result 2] compared to the strongest
prior baseline \cite{PLACEHOLDER_baseline}.

\textbf{Contributions.} This paper makes the following contributions:

\begin{itemize}
    \item We identify [a specific limitation of prior work] and characterize
          it on [N] benchmark designs (Section~\ref{sec:background}).
    \item We propose \method{}, a [class of method] that addresses this
          limitation by [the key idea] (Section~\ref{sec:method}).
    \item We provide [analytical / theoretical] evidence that \method{}
          [property], with complexity [\(O(n \log n)\) or similar]
          (Section~\ref{sec:method}).
    \item We evaluate \method{} on the [benchmark suite] and report
          [headline result] over [N] baselines, including [B1] and [B2]
          (Section~\ref{sec:experiments}).
    \item We release the implementation as open source at
          \texttt{[URL or anonymous mirror during review]}
          (Section~\ref{sec:experiments}).
\end{itemize}

The rest of this paper is organized as follows.
Section~\ref{sec:background} reviews relevant background and prior work.
Section~\ref{sec:method} presents \method{} in detail.
Section~\ref{sec:experiments} describes the experimental setup and
results. Section~\ref{sec:discussion} discusses limitations and threats
to validity. Section~\ref{sec:conclusion} concludes.
